% GENERATED FILE. Edit canonical lessons, then run npm run generate:books.
% canonical-source-hash: fnv1a64:689bfaeb4e20edfc
% canonical-lessons: SA-C115-maksika, SA-C115-masakah, SA-C115-pipilika, SA-C115-citrapatangah, SA-C115-kukkutah

\chapter{Fly, Mosquito, Ant, Butterfly, Rooster}
\label{ch:sa-fly-mosquito-ant-butterfly-rooster}
\hlchaptermodality{115}

\begin{chapteropening}
\textbf{By the end of this chapter:} \emph{I can say a fly, a mosquito, an ant, a butterfly and a rooster.}

Complete the last lesson of chapter 115: i can say a fly, a mosquito, an ant, a butterfly and a rooster.
\end{chapteropening}

\section[makṣikā]{\sk{मक्षिका} (makṣikā) — a fly}

\label{lesson:SA-C115-maksika}



\begin{quote}
Before the new one: say the Sanskrit for a heron, then the Sanskrit for a bee.
\end{quote}

\subsection*{You'll want to know: \sk{मक्षिका}}
\textbf{\sk{मक्षिका}} — \emph{makṣikā} — \textquotedblleft{}a fly\textquotedblright{}.

\begin{grammarlens}[title={What you've built}]
One more word for this chapter.
\end{grammarlens}

\subsection*{Your turn}
\begin{itemize}
\raggedright
  \item \emph{Say it:} \emph{makṣikā}
  \item \emph{Say it:} \emph{makṣikā}, once more
  \item \emph{Say it:} say \emph{makṣikā}
  \item \emph{Recall:} say the Sanskrit for a heron, then the Sanskrit for a bee, then say \emph{makṣikā} again
\end{itemize}

\begin{tcolorbox}[breakable,colback=teal!4,colframe=teal!35!black,title={Before you move on}]
What does \sk{मक्षिका} mean? (\textquotedblleft{}A fly\textquotedblright{}.) Say it once more.
\end{tcolorbox}
\section[maśakaḥ]{\sk{मशकः} (maśakaḥ) — a mosquito}

\label{lesson:SA-C115-masakah}



\begin{quote}
Before the new one: say the Sanskrit for a bee, then the Sanskrit for a fly.
\end{quote}

\subsection*{You'll want to know: \sk{मशकः}}
\textbf{\sk{मशकः}} — \emph{maśakaḥ} — \textquotedblleft{}a mosquito\textquotedblright{}.

\begin{grammarlens}[title={What you've built}]
One more word for this chapter.
\end{grammarlens}

\subsection*{Your turn}
\begin{itemize}
\raggedright
  \item \emph{Say it:} \emph{maśakaḥ}
  \item \emph{Say it:} \emph{maśakaḥ}, once more
  \item \emph{Say it:} say \emph{maśakaḥ}
  \item \emph{Recall:} say the Sanskrit for a bee, then the Sanskrit for a fly, then say \emph{maśakaḥ} again
\end{itemize}

\begin{tcolorbox}[breakable,colback=teal!4,colframe=teal!35!black,title={Before you move on}]
What does \sk{मशकः} mean? (\textquotedblleft{}A mosquito\textquotedblright{}.) Say it once more.
\end{tcolorbox}
\section[pipīlikā]{\sk{पिपीलिका} (pipīlikā) — an ant}

\label{lesson:SA-C115-pipilika}



\begin{quote}
Before the new one: say the Sanskrit for a fly, then the Sanskrit for a mosquito.
\end{quote}

\subsection*{You'll want to know: \sk{पिपीलिका}}
\textbf{\sk{पिपीलिका}} — \emph{pipīlikā} — \textquotedblleft{}an ant\textquotedblright{}.

\begin{grammarlens}[title={What you've built}]
One more word for this chapter.
\end{grammarlens}

\subsection*{Your turn}
\begin{itemize}
\raggedright
  \item \emph{Say it:} \emph{pipīlikā}
  \item \emph{Say it:} \emph{pipīlikā}, once more
  \item \emph{Say it:} say \emph{pipīlikā}
  \item \emph{Recall:} say the Sanskrit for a fly, then the Sanskrit for a mosquito, then say \emph{pipīlikā} again
\end{itemize}

\begin{tcolorbox}[breakable,colback=teal!4,colframe=teal!35!black,title={Before you move on}]
What does \sk{पिपीलिका} mean? (\textquotedblleft{}An ant\textquotedblright{}.) Say it once more.
\end{tcolorbox}
\section[citrapataṅgaḥ]{\sk{चित्रपतंगः} (citrapataṅgaḥ) — a butterfly}

\label{lesson:SA-C115-citrapatangah}



\begin{quote}
Before the new one: say the Sanskrit for a mosquito, then the Sanskrit for an ant.
\end{quote}

\subsection*{You'll want to know: \sk{चित्रपतंगः}}
\textbf{\sk{चित्रपतंगः}} — \emph{citrapataṅgaḥ} — \textquotedblleft{}a butterfly\textquotedblright{}.

\begin{grammarlens}[title={What you've built}]
One more word for this chapter.
\end{grammarlens}

\subsection*{Your turn}
\begin{itemize}
\raggedright
  \item \emph{Say it:} \emph{citrapataṅgaḥ}
  \item \emph{Say it:} \emph{citrapataṅgaḥ}, once more
  \item \emph{Say it:} say \emph{citrapataṅgaḥ}
  \item \emph{Recall:} say the Sanskrit for a mosquito, then the Sanskrit for an ant, then say \emph{citrapataṅgaḥ} again
\end{itemize}

\begin{tcolorbox}[breakable,colback=teal!4,colframe=teal!35!black,title={Before you move on}]
What does \sk{चित्रपतंगः} mean? (\textquotedblleft{}A butterfly\textquotedblright{}.) Say it once more.
\end{tcolorbox}
\section[kukkuṭaḥ]{\sk{कुक्कुटः} (kukkuṭaḥ) — a rooster}

\label{lesson:SA-C115-kukkutah}



\begin{quote}
Before the new one: say the Sanskrit for an ant, then the Sanskrit for a butterfly.
\end{quote}

\subsection*{You'll want to know: \sk{कुक्कुटः}}
\textbf{\sk{कुक्कुटः}} — \emph{kukkuṭaḥ} — \textquotedblleft{}a rooster\textquotedblright{}.

\begin{grammarlens}[title={What you've built}]
One more word for this chapter.
\end{grammarlens}

\subsection*{Your turn}
\begin{itemize}
\raggedright
  \item \emph{Say it:} \emph{kukkuṭaḥ}
  \item \emph{Say it:} \emph{kukkuṭaḥ}, once more
  \item \emph{Say it:} say \emph{kukkuṭaḥ}
  \item \emph{Recall:} say the Sanskrit for an ant, then the Sanskrit for a butterfly, then say \emph{kukkuṭaḥ} again
\end{itemize}

\begin{tcolorbox}[breakable,colback=teal!4,colframe=teal!35!black,title={Before you move on}]
What does \sk{कुक्कुटः} mean? (\textquotedblleft{}A rooster\textquotedblright{}.) Say it once more.
\end{tcolorbox}
